% ============================================================
% Front Matter: Preface, Acknowledgments, How to Use This Book
% ============================================================

\chapter*{Preface}
\addcontentsline{toc}{chapter}{Preface}

On the morning of August 8, 1900, David Hilbert stood before the assembled mathematicians of the International Congress of Mathematicians in Paris and delivered a lecture that would change the course of mathematics. In that lecture, Hilbert proposed twenty-three problems---open questions spanning the breadth of mathematics---that he believed should guide research in the coming century.

More than a century later, the legacy of that lecture endures. Hilbert's problems have served as beacons for mathematical research, inspiring entire fields and leading to some of the deepest theorems of modern mathematics. Some have been solved, some remain open, and some have evolved into questions that Hilbert could never have imagined.

This book is written for undergraduate students who wish to understand what Hilbert asked, why it mattered, and where mathematics stands today in addressing his challenges. No graduate-level background is assumed. The reader should be comfortable with calculus, basic linear algebra, and some exposure to proofs, but the mathematical details are kept at a level accessible to a motivated undergraduate.

Each chapter addresses one of Hilbert's problems. We begin with historical context---what was known in 1900 and what Hilbert hoped for. We then explain the mathematical content of the problem in modern language, describe the key ideas that led to its resolution (or progress toward it), and state the current status. Throughout, we emphasize the \emph{ideas} over technical machinery, aiming to convey the beauty and depth of these mathematical questions.

Mathematics is not a spectator sport. At the end of each chapter, you will find exercises. Some are computational, some are conceptual, and some ask you to explore further on your own. Working through them will deepen your understanding and---we hope---inspire you to learn more about the remarkable areas of mathematics that Hilbert's problems opened up.

\vspace{2em}

\noindent\textbf{Prerequisite guide by chapter.}
The book's baseline requirement is calculus, basic linear algebra, and mathematical maturity (proofs). However, the problems span the entire mathematical curriculum, and some chapters demand more. The table below maps each problem to the additional courses that will help you get the most out of that chapter.

\begin{center}
\small
\renewcommand{\arraystretch}{1.2}
\begin{tabular}{ccl}
\toprule
\textbf{Problem} & \textbf{Title} & \textbf{Recommended preparation} \\
\midrule
1 & Continuum Hypothesis & --- \\
2 & Consistency of Arithmetic & --- \\
3 & Equality of Volumes & --- \\
4 & Shortest Path = Straight Line & --- \\
5 & Continuous Groups & Linear algebra, group theory \\
6 & Axiomatization of Physics & --- \\
7 & Irrationality of $e^\pi$ & --- \\
8 & Goldbach Conjecture & --- \\
9 & Reciprocity Laws & Abstract algebra (Galois theory) \\
10 & Decision for Diophantine Eqns & --- (or computability theory) \\
11 & Quadratic Forms & Abstract algebra (rings, fields) \\
12 & Kronecker's Theorem & Abstract algebra (Galois theory) \\
13 & Solving 7th Degree Eqns & Abstract algebra (Galois theory) \\
14 & Constructive Invariants & Abstract algebra (rings, modules) \\
15 & Schubert's Enumerative Calc. & --- \\
16 & Topology of Varieties & Algebraic topology \\
17 & Representation by Squares & --- \\
18 & Tessellation of Space & --- \\
19 & Regularity of PDE Solutions & Real analysis (PDEs) \\
20 & Boundary Value Problems & Real analysis (PDEs) \\
21 & Monodromy Problem & Complex analysis \\
22 & Uniformization & Complex analysis \\
23 & Calculus of Variations & Real analysis \\
\bottomrule
\end{tabular}
\end{center}

Chapters marked ``---'' require only the baseline. Each chapter defines its own concepts, so you can always \emph{read} a chapter above your preparation; starred exercises ($\star$) in advanced chapters will be harder without the listed background. The book is designed to be read non-linearly---skip around freely.

\chapter*{Acknowledgments}

\chapter*{Acknowledgments}
\addcontentsline{toc}{chapter}{Acknowledgments}

This book draws on over a century of mathematical scholarship. We are indebted to the many mathematicians who solved, advanced, and illuminated Hilbert's problems. Particular gratitude is due to Jeremy Gray, whose historical scholarship on Hilbert's problems has been invaluable; to Constance Reid, whose biography of Hilbert remains the definitive account; and to the many mathematicians who have written expositions of individual problems, making them accessible to broader audiences.

We also thank the generations of students whose questions and curiosity helped shape this exposition. Any errors that remain are, of course, our own.

\chapter*{How to Use This Book}
\addcontentsline{toc}{chapter}{How to Use This Book}

\begin{itemize}
    \item \textbf{Each chapter} addresses one of Hilbert's 23 problems and is self-contained. You may read them in order or jump to problems that interest you.

    \item \textbf{Historical context} appears at the beginning of each chapter, explaining what was known in 1900.

    \item \textbf{The problem} is stated in modern mathematical language, as close to Hilbert's original intent as possible.

    \item \textbf{Key developments} describe the major results and ideas that have addressed the problem.

    \item \textbf{Current status} summarizes where the problem stands today, using the following conventions:
    \begin{itemize}
        \item \textbf{Solved:} A complete answer (or proof of impossibility) is known.
        \item \textbf{Partially solved:} Significant progress has been made, but the full problem remains open.
        \item \textbf{Unsolved:} Little progress, or the problem remains wide open.
        \item \textbf{Independent:} The problem cannot be decided within standard axiom systems (a different kind of resolution).
    \end{itemize}

    \item \textbf{Exercises} range from routine to challenging. star-marked exercises ($\star$) are more difficult.

    \item \textbf{Appendices} provide a timeline, glossary of terms, and suggestions for further reading.
\end{itemize}
